\section{Idea、实验与时间规划}

\subsection{一条常见的阶段时间表}
\begin{center}
\begin{tabular}{clp{7.2cm}}
\toprule
阶段 & 约需时间 & 任务 \\
\midrule
想 idea & 1--2 个月 & 规划任务、发现 technical challenge / failure case、构思解法 \\
做实验 & 2--3 个月 & 简单实验验证 $\rightarrow$ 简单数据改进 $\rightarrow$ 真实数据调通 \\
写论文 & 约 1 个月 & 按投稿节点倒排、写作、自查、修改 \\
\bottomrule
\end{tabular}
\end{center}
实际往往更长；假期与课余集中时间对前期积累尤其重要。

\subsection{两种科研推进方式}
\paragraph{Idea-driven（跟进式）}
看到有趣的 SOTA，发现 failure case，提出新算法修补。\\
优点：上手快，适合新手。\\
缺点：影响力依赖提升幅度；易撞车；创新空间受原任务束缚。

\paragraph{Goal-driven（目标牵引）}
先设定长期目标与 roadmap，再选任务、挖 failure case、设计解法。\\
优点：更易做出成体系、有动机的工作。\\
缺点：规划 roadmap 本身很难，需要导师指导与领域积累。

\subsection{如何发现与解决 Failure Cases}
\begin{enumerate}
  \item 探索算法上限，在更难 case / 新数据设定上找失败模式；
  \item 把表面现象抽象成真正的 technical challenge；
  \item 从自己的「武器库」（可复用的 technique / insight）中组合解法；
  \item 好的工作往往不是简单 A+B，而是把 A 与 B 巧妙结合。
\end{enumerate}
当领域出现「新锤子」（如扩散模型、大模型范式）时，用它去推进自己 roadmap 上的里程碑，常更容易做出有影响力的工作。

\subsection{实验策略：先简单后真实}
\textbf{简单实验}：尽量只保留你要验证的那个 core challenge，排除无关干扰，先证明解法「在该挑战上 work」。\\
然后再到真实数据上把系统调通，并补齐对比与消融。

\subsection{实验记录最低要求}
日期、步骤、想法、失败实验、配置、结果图表、下一步计划，都应可回溯。实验记录是成果的原始证据。

\subsection{经典模型精读清单（代码与论文）}
建议结合李沐论文精读仓库：\url{https://github.com/mli/paper-reading}\\
优先顺序示例：ResNet $\rightarrow$ Transformer $\rightarrow$ GAN $\rightarrow$ 对比学习 $\rightarrow$ Swin Transformer $\rightarrow$ CLIP $\rightarrow$ Two-Stream $\rightarrow$ 多模态学习。
